\chapter{Att kontrollera ett påstående}
\label{app:verifiera}

Det här materialet gör, som allt undervisningsmaterial, en rad
påståenden.  De flesta går att pröva vid datorn: att ett attribut med
dubbla understreck inte går att nå utifrån visar materialet med en
körning, och du kan göra om den själv.  Men några av påståendena
handlar om vad forskningen säger, och dem kan ingen körning avgöra:
\begin{enumerate}
  \item att nybörjare systematiskt missförstår vad en klass och ett
    objekt är --- att klassen ses som en samling av sina objekt och att
    två objekt inte kan ha samma attributvärden
    (\cref{sec:med-klasser,sec:klasser-objekt,sec:dundermetoder});
  \item att inkapsling är en etablerad princip i objektorienterad
    programmering (\cref{sec:tva-exempel});
  \item att man bör föredra komposition framför arv, ett etablerat råd
    från Gamma, Helm, Johnson och Vlissides (\cref{sec:arv});
  \item att en klass med ett enda ansvar följer principen \emph{SRP},
    som är Robert C. Martins (\cref{sec:tva-exempel});
  \item att det är bra att inte upprepa sig i ett program --- principen
    \emph{DRY} --- och att den kommer från Hunt och Thomas
    (\cref{sec:arv}).
\end{enumerate}
Hur vet man om sådant stämmer?  Samma fråga möter dig varje gång du läser
en lärobok, en webbsida eller ett svar från en språkmodell --- och varje
gång du själv ska skriva en rapport.  Den här bilagan sammanfattar
arbetsgången, bilagorna som följer visar den genomförd och
\cref{tab:paastaenden} sammanfattar frågorna och svaren; SRP och DRY är
prövade i andra föreläsningar.  Arbetsgången beskrivs utförligare i bilagan med samma namn i föreläsningen
\emph{Algoritmiskt tänkande}.

\begin{table}[htbp]
  \centering
  \caption{Påståendena som frågor, och var svaren finns.}
  \label{tab:paastaenden}
  \begin{tabular}{@{}lp{0.52\linewidth}p{0.3\linewidth}@{}}
    \toprule
    Bilaga & Fråga & Svar \\
    \midrule
    \Cref{app:missuppfattningar} & Missförstår nybörjare vad en klass
      och ett objekt är, och i så fall hur? & Ja, med förbehåll \\
    \Cref{app:inkapsling} & Är inkapsling en etablerad princip i
      objektorienterad programmering? & Ja, med förbehåll \\
    \Cref{app:komposition} & Bör man föredra komposition framför arv,
      och kommer rådet från Gamma med flera? & Ja, med förbehåll \\
    --- & Är SRP Robert C. Martins princip, och är den etablerad? &
      Redan belagt i \emph{Funktioner}, bilaga~B \\
    --- & Är det bra att inte upprepa sig, och kommer DRY från Hunt och
      Thomas? & Redan belagt i \emph{Algoritmiskt tänkande}, bilaga~E \\
    \bottomrule
  \end{tabular}
\end{table}

\section{Gör påståendet till en fråga}

Ett påstående går att kontrollera först när det är tydligt vad som skulle
kunna visa att det är \emph{fel}.  Formulera det därför som en fråga med
ett svar.  \enquote{Missförstår nybörjare vad en klass är?} kan besvaras
med ja eller nej; \enquote{klasser är viktiga} kan det inte, eftersom
\enquote{viktiga} är ett omdöme och inte ett faktum.  Ofta döljer sig
flera påståenden i ett: \enquote{en etablerad princip} är två frågor ---
är den en \emph{princip}, och är den \emph{etablerad} som sådan --- och de
kräver olika slags källor.

\section{Välj rätt sorts källa}

Olika frågor besvaras av olika källor.
\begin{description}
  \item[Uppslagsverk] för vad ett ord \emph{betyder}.  För vad ett
    programspråk \emph{gör} är språkets egen dokumentation motsvarande
    källa; därför står Pythons språkreferens i fotnoterna till
    \cref{sec:dundermetoder}, med adress och datum.
  \item[Primärkällan] för var en idé \emph{kommer ifrån}, och för ett
    specifikt resultat.  Vill man veta vad studenter faktiskt trodde
    läser man studien som frågade dem, inte en lärobok som återger den.
  \item[Översikter] (\foreignlanguage{english}{reviews}) för om något är
    \emph{etablerat}, eller för vad forskningen \emph{sammantaget}
    säger.
  \item[Enskilda studier] för ett specifikt resultat --- med förbehållet
    att en enskild studie bara visar att något observerats en gång, i en
    viss population.
\end{description}
Var man hittar dem: universitetsbibliotekets söktjänst, och databaser som
Scopus, Web of Science, IEEE Xplore, DBLP (datalogi) och OpenAlex
(öppen).  I bilagorna användes kommandoradsverktyget
\texttt{scholar}\footnote{Paketet \texttt{scholarcli} på PyPI:
  \url{https://pypi.org/project/scholarcli/}.}, som ställer samma fråga
till flera databaser på en gång och sparar allt som hände i en
\emph{session}.  Det du får ut av bibliotekets söktjänst är detsamma,
bara långsammare.

\section{Sök så att du kan ha fel}

Den som söker på namn hon redan känner till hittar det hon väntar sig.
Sök därför \emph{ämnesdrivet} --- på begrepp, inte författare.  Alla
sökningar i bilagorna är gjorda så.  Några praktiska regler:
\begin{itemize}
  \item Håll frågorna korta.  Flera korta frågor slår en lång.
  \item Ställ \emph{samma} frågor i alla databaser.  Databaserna har
    olika frågespråk, och en fråga som fungerar i en kan ge noll träffar
    i en annan av rena tekniska skäl; det är inte ett svar utan en fråga
    som ställdes fel, och den ska ställas om.
  \item Sök också efter \emph{invändningar}: lägg till ord som
    \enquote{critique}, \enquote{limitations} eller \enquote{harmful}.
    Ett påstående som överlevt en motsökning är värt mer än ett som
    aldrig prövats, och en invändning som hittas blir ett
    \emph{förbehåll} i svaret, inte ett nederlag.
  \item Räkna med att databasernas rankning missar somligt.  Då hämtas
    källan som \emph{känt dokument}: man vet att den finns och slår upp
    den direkt.  Det är i sin ordning, bara det sägs.
\end{itemize}

\section{Läs i fulltext, och skriv ner hur}

En träff i en databas är inte ett belägg.  Sammanfattningen
(\foreignlanguage{english}{abstract}) räcker inte heller: den säger vad
författarna vill ha sagt, inte vad de visat.  Läs hela texten och leta
upp \emph{det ställe} som styrker påståendet --- ett ordagrant citat med
sidnummer.  Hittar du inget sådant ställe stöder källan inte påståendet,
hur passande titeln än är.  Räkna aldrig en källa du inte läst.  Går
fulltexten inte att nå --- så är det med en av källorna i
\cref{app:missuppfattningar} och de flesta i \cref{app:komposition}
--- skriver man det, och nöjer sig med det sammanfattningen säger.

Skriv sedan ner hur du gick till väga, så att någon annan kan följa
spåret.  I det här materialet står den anteckningen som en kommentar
ovanför varje post i litteraturlistans källfil, med fasta fält:
\begin{description}
  \item[\texttt{CLAIM}] vilket påstående källan ska styrka;
  \item[\texttt{FOUND-VIA}] hur den hittades --- sökfråga och databas,
    eller \enquote{känt dokument};
  \item[\texttt{PICKED}] varför just den, framför andra träffar;
  \item[\texttt{QUOTE}] det ordagranna citatet, med sida;
  \item[\texttt{VERIFIED}] att fulltexten lästs, och var;
  \item[\texttt{COUNTER}] motsökningen och vad den gav;
  \item[\texttt{DATE}] när.
\end{description}
Det tar ett par minuter per källa, och det är det enda som gör det
möjligt att om ett år svara på frågan \enquote{hur vet du det?}

\section{Dra slutsatsen --- med förbehåll}

Svaret skrivs ut tillsammans med det som talar emot.  Bilagorna svarar
alla \enquote{ja, med förbehåll} (\cref{tab:paastaenden}).
Förbehållen är inte svagheter i svaret.  De \emph{är} svaret, och de
avgör hur påståendet används i materialet.  Varje kapitel slutar därför
med ett avsnitt \emph{Fortsatt arbete}: vad just den sökningen lämnade
ogjort, och vad som skulle behöva undersökas för ett säkrare svar.

\section{En anmärkning om språkmodeller}

Sökningarna i bilagorna gav hundratals träffar var.  En språkmodell
användes för att \emph{sortera} dem --- hör träffen till ämnet över
huvud taget, och åt vilket håll pekar den? --- så att gallringen kunde
granskas.  Modellen sorterade, men varje källa som citeras ska läsas och
väljas av en människa, och även sorteringen ska bekräftas av en
människa.  Modellens klassning står därför med sin konfidens i
tabellerna, så att den kan ifrågasättas.  Använd gärna språkmodeller för
att hitta och sortera; låt dem aldrig vara det sista ledet.

\chapter{Missförstår nybörjare klasser och objekt?}
\label{app:missuppfattningar}
\chapterprecis{Författaren har ännu inte granskat resultaten i den här
  bilagan i sin helhet.}

I \cref{sec:med-klasser,sec:klasser-objekt,sec:dundermetoder} ställer
materialet tre frågor som inte fanns i tidigare versioner: var de två
objekten finns när de skapats från samma klass, vad som händer med det
ena objektet när det andra byter namn, och vad \mintinline{python}{==}
ger för två objekt med samma attributvärden.  Frågorna finns för att tre
missuppfattningar om klasser och objekt är belagda hos nybörjare: att
klassen är en samling av sina objekt snarare än en mall för att skapa
dem\autocite{Ragonis2005OOP}, att ett objekts attribut är dess identitet
så att två objekt inte kan ha samma värden\autocite{Holland1997}, och
mer allmänt att studenter saknar en grundläggande förståelse av
objektbegreppet\autocite{Kaczmarczyk2010}.  Frågan som ska besvaras är:
\emph{missförstår nybörjare vad en klass och ett objekt är, och i så
fall hur?}  \Cref{sec:miss-metod} beskriver hur påståendet prövats,
\cref{sec:miss-resultat} vad som fanns, \cref{sec:miss-diskussion} vad
som begränsar svaret, \cref{sec:miss-slutsats} ger det och
\cref{sec:miss-fortsatt} säger vad som återstår att göra;
\cref{tab:sok-missuppfattningar} sist i kapitlet listar de träffar som
rör påståendet.

\section{Metod}
\label{sec:miss-metod}

Påståendet är ett av dem som författarens egen genomgång av
missuppfattningar i programmering \parencite{SooriBosk2026} redan har
prövat: Ragonis och
Ben-Ari, Holland med flera och Kaczmarczyk med flera kommer alla tre
från den genomgångens litteraturprotokoll, där de hittades i en
manuell sökning och lästes.  Här gjordes två saker till.  Källorna
kontrollerades på nytt mot fulltexten, med olika utfall:
Kaczmarczyk med fleras artikel lästes i författarens egen öppna kopia
och tabellen med de fyra temana finns där ordagrant; Holland med fleras
artikel, som i genomgången bara kunde kontrolleras mot sammanfattningen,
lästes här i författarnas manuskript i Open Universitys arkiv, så att
de enskilda missuppfattningarna kunde citeras ordagrant; Ragonis och
Ben-Aris artikel är stängd och kunde inte nås, så det citat som
materialet bygger på är ärvt från genomgången och står som sådant i
källfilen.

Därtill söktes frågan ämnesdrivet och tvåsidigt den 3 september 2026
med verktyget \texttt{scholar}, mot de databaser som svarade: OpenAlex,
Web of Science och Scopus, samt DBLP, som svarade på en del av frågorna
och avvisade andra med ett serverfel.  IEEE Xplore hade nått sin
dagsgräns och Semantic Scholar avvisade sin nyckel; båda saknas.
Stödfrågorna söker missuppfattningar och begreppsliga svårigheter hos
nybörjare i objektorienterad programmering, dels allmänt, dels om just
klass, objekt, instans och identitet, dels som översikter; motfrågorna
söker med samma begrepp efter studier där svårigheterna är små, saknas
eller beror på undervisningsordningen, och efter replikationer och
motsägelser.  \Cref{tab:sok-queries-B} visar frågorna och antalet
träffar per databas.

% Author's note (verification and reproduction): see the block above the
% hit-list \input at the end of this chapter.  The query table is
% transcribed from the session's query records (scholar sessions show
% classes-missuppfattningar) and the per-provider hit counts recorded
% there.
\begin{table}[!htb]
  \centering
  \caption{Sökfrågor och antal träffar per databas,
    missuppfattningar om klasser och objekt, 3 september 2026, verktyget
    \texttt{scholar}.  S = stödfråga, M = motfråga.  Databaser: OpenAlex
    (OA; inte \foreignlanguage{english}{open access}), DBLP, Web of
    Science (WoS), Scopus.  Högst 40 träffar hämtades per databas och
    fråga, så 40 betyder \enquote{minst 40}.  Frågan visas i den
    booleska form som skickades till OpenAlex; Web of Science fick samma
    sträng som \texttt{TS=(\dots)}, Scopus som
    \texttt{TITLE-ABS-KEY(\dots)} och DBLP en kort ordlista med samma
    begrepp (DBLP indexerar bara titlar och kräver att varje ord
    förekommer; ordlistorna står i författarens anteckning).  IEEE Xplore
    hade nått sin dagsgräns och Semantic Scholar avvisade sin nyckel; de
    saknas.}
  \label{tab:sok-queries-B}
  \footnotesize
  \begin{tabular}{@{}lp{0.55\linewidth}rrrr@{}}
    \toprule
    & Nyckelord & OA & DBLP & WoS & Scopus \\
    \midrule
    S1 & (\enquote{object-oriented} \textsc{or} \enquote{object oriented} \textsc{or} OOP) \textsc{and} (novice \textsc{or} novices \textsc{or} beginners \textsc{or} CS1) \textsc{and} (misconceptions \textsc{or} \enquote{conceptual difficulties} \textsc{or} \enquote{mental models})
       & 40 & 4 & 25 & 38 \\
    S2 & (class \textsc{or} classes \textsc{or} object \textsc{or} objects \textsc{or} instantiation) \textsc{and} (template \textsc{or} instance \textsc{or} identity \textsc{or} \enquote{same object}) \textsc{and} (novice \textsc{or} novices \textsc{or} students) \textsc{and} (misconception \textsc{or} misconceptions \textsc{or} difficulties \textsc{or} conceptions)
       & 40 & 9 & 39 & 40 \\
    S3 & (\enquote{object-oriented programming} \textsc{or} \enquote{object-oriented}) \textsc{and} (learning \textsc{or} teaching \textsc{or} education) \textsc{and} (misconceptions \textsc{or} difficulties) \textsc{and} (review \textsc{or} survey \textsc{or} systematic)
       & 40 & 1 & 36 & 40 \\
    \midrule
    M1 & (\enquote{object-oriented} \textsc{or} \enquote{object oriented} \textsc{or} \enquote{objects-first} \textsc{or} \enquote{objects first}) \textsc{and} (novice \textsc{or} novices \textsc{or} beginners \textsc{or} CS1) \textsc{and} (easier \textsc{or} \enquote{no difference} \textsc{or} \enquote{not more difficult} \textsc{or} \enquote{no significant} \textsc{or} critique \textsc{or} criticism \textsc{or} overstated)
       & 40 & 0 & 36 & 40 \\
    M2 & (\enquote{object-oriented} \textsc{or} OOP) \textsc{and} (novices \textsc{or} students) \textsc{and} (misconceptions \textsc{or} difficulties) \textsc{and} (replication \textsc{or} contradictory \textsc{or} \enquote{not supported} \textsc{or} limitations \textsc{or} \enquote{do not})
       & 40 & 1 & 15 & 23 \\
    \bottomrule
  \end{tabular}
\end{table}

Sammanlagt lämnade databaserna 415 poster.  Flera databaser lämnar samma
arbete, och två poster räknas som ett arbete när titlarna är desamma
bortsett från versaler och skiljetecken; efter den sammanslagningen
återstår 387 arbeten, och det är arbeten som räknas här och i
\cref{tab:sok-missuppfattningar}.  Alla klassades: först med
en språkmodell mot en beskrivning av frågan och fyra kategorier ---
stöder påståendet, kvalificerar eller motsäger det, angränsande delämne,
annat ämne --- och sedan för hand: varje post som modellen satt i den
andra kategorin, och varje post med konfidens under 0,7, lästes i
sammanfattning.  De källor som kapitlet citerar lästes i fulltext där
den gick att nå; för de övriga står det i källfilen att bara
sammanfattningen lästs, och \cref{sec:miss-diskussion} återkommer till
det.  Hur varje
källa hittades och verifierades är antecknat på det sätt
\cref{app:verifiera} beskriver.

\Cref{tab:sok-missuppfattningar} rättades för hand vid genomläsningen,
och rättelserna redovisas därför att en rättelse man inte kan
kontrollera inte är någon rättelse.  Fem titlar bar spår av
\textsc{html}-tecken i leverantörernas titlar --- citattecknen kring
\enquote{object}, \enquote{class} och \enquote{this} och apostroferna i
\enquote{Students'} hade kommit ut som obearbetade styrsekvenser eller
som svenska babel-genvägar --- och en sjätte var avhuggen mitt i titeln
av exportens teckengräns.  Två arbeten stod på tre respektive två rader,
ett per databas, och står nu på var sin rad med databaserna uppräknade.
Och en rad bytte kategori: Holland med fleras artikel stod som stödjande
med modellens konfidens, men den citeras i kapitlet och är läst i
fulltext, så den står nu som citerad.

\section{Resultat}
\label{sec:miss-resultat}

\paragraph{Klassen som samling.}  \Textcite{Ragonis2005OOP} följde
gymnasieelever som lärde sig Java med BlueJ under två läsår och
katalogiserade 58 uppfattningar och svårigheter i fyra kategorier, varav
den första gäller just klass mot objekt.  Bland dem står, enligt
genomgångens läsning av artikeln, att en klass är en samling av objekt
snarare än en mall för att skapa dem.  Två senare studier av samma två
begrepp bekräftar bilden.  \Textcite{EckerdalThune2005} undersökte
förstaårsstudenters uppfattningar av \enquote{objekt} och \enquote{klass}
fenomenografiskt och använde variationsteori för att
\textcquote{EckerdalThune2005}{pin-point what the students need to be
able to discern in order to gain a \mkbibquote{rich} understanding of
these concepts} --- samma verktyg som det här materialet är byggt med.
Deras kategorier (s.~91) är inkluderande nivåer av en riktig förståelse
--- objektet som \enquote{a piece of code}, som \enquote{something that
is active in the program}, som \enquote{a model of some real world
phenomenon} --- och artikeln studerar uttryckligen uppfattningar, inte
missuppfattningar (s.~89); det den bekräftar är att skillnaden mellan
klass och objekt är svår att nå.  Dess råd är att alltid arbeta med flera
instanser av varje klass (s.~92--93), vilket det här materialet gör.
\Textcite{Xinogalos2015} ger tio år senare frekvenser:
\textcquote{Xinogalos2015}{[n]early half the students seem to comprehend
the modeling and static/dynamic aspects of the concepts
\mkbibquote{object} and \mkbibquote{class}} --- och alltså nästan hälften
inte.

\paragraph{Attributen som identitet.}  \Textcite{Holland1997} beskriver,
ur sin undervisning i Smalltalk, hur det klassiska bankkontoexemplet med
ett namnattribut får studenter att blanda ihop attributet med objektets
identitet, och räknar upp följdmissuppfattningarna:
\textcquote{Holland1997}{if you have two different variables, they must
refer to two different objects; [\dots] two objects of the same class
with the same state are the same object; two objects with the same value
for the name attribute are the same object}.  Deras botemedel är just
det materialet gör i \cref{sec:dundermetoder}: \textcquote{Holland1997}{
[r]ather than try to deal with these misconceptions by arguing or talking
about them, the easiest approach is to immediately let the students
experiment with a set of counter examples}.  Samma artikel varnar för
exempel där objekt bara är passiva poster ---
\textcquote{Holland1997}{students may come to tacitly assume that all
objects are simple, inert records} --- vilket är skälet till att
personklassen i \cref{sec:tva-exempel} har metoder som gör något, inte
bara läser attribut.

\paragraph{Objektbegreppet i stort.}  I intervjuer med studenter i en
första programmeringskurs fann \textcite{Kaczmarczyk2010} fyra teman av
missuppfattningar, varav det tredje lyder ordagrant:
\textcquote{Kaczmarczyk2010}{T3: Students lack a basic understanding of
the Object concept}.  \Textcite{SandersThomas2007} fann i studenternas
egna program i en objektförst-kurs \textcquote{SandersThomas2007}{concrete
evidence of students learning these concepts while also displaying some
of these misconceptions}.  Sökningen gav därutöver nio arbeten som
stöder påståendet utan att ha lästs i fulltext; de står som stödjande i
\cref{tab:sok-missuppfattningar}.  Holland med fleras artikel kom också
upp i sökningen, men den är läst i fulltext och citeras här, så den står
som citerad i tabellen.

\paragraph{Motsökningen.}  Ingen studie som hävdar att svårigheterna
saknas kom fram.  Det som kom fram kvalificerar: svårigheten beror på
omständigheterna.  \Textcite{Wiedenbeck1999} jämförde nybörjare som lärt
sig objektorienterad respektive procedurell programmering och fann
\textcquote{Wiedenbeck1999}{no significant difference} för korta program
men sämre förståelse hos de objektorienterade för ett längre, vilket de
knyter till \textcquote{Wiedenbeck1999}{a longer learning curve};
\textcite{KunkleAllen2016} fann \textcquote{KunkleAllen2016}{significant
differences in student performance based on instructional language and
approach}; och en jämförelse av objektförst mot procedurellt-först i
\cref{tab:sok-missuppfattningar} fann att objektförst gav bättre
förståelse.  Svårigheten är alltså verklig men inte oberoende av hur
och i vilken ordning man undervisar.

\section{Diskussion och begränsningar}
\label{sec:miss-diskussion}

Tre saker begränsar svaret.  Källorna gäller andra språk än Python:
Java hos Ragonis och Ben-Ari, Eckerdal och Thuné, Xinogalos, Kaczmarczyk
med flera och Sanders och Thomas, Smalltalk hos Holland med flera.
Missuppfattningarna handlar om begreppen klass och objekt, inte om
syntax, så överföringen är rimlig, men hur vanliga de är hos just våra
studenter säger studierna inget om.  Holland med fleras artikel är
dessutom en erfarenhetsrapport från undervisning, inte en mätning.
Verifieringen är ojämn: Holland med flera och Kaczmarczyk med flera
lästes i fulltext, liksom Eckerdal och Thunés artikel i författarens
eget exemplar, men Ragonis och Ben-Aris artikel är stängd och citatet
ärvt från genomgången, och de två övriga bekräftande studierna liksom de
två kvalificerande kunde bara läsas i sammanfattning, eftersom ACM:s
digitala bibliotek avvisar automatisk hämtning; det står i källfilen och
bör kontrolleras av den som har tillgång.  Slutligen är täckningen
ofullständig: IEEE Xplore och Semantic Scholar saknas, DBLP svarade
ojämnt, träffantalen i \cref{tab:sok-queries-B} är tak, och klassningen
är modellstödd, så antalet stödjande poster är en övre gräns.

\section{Slutsats}
\label{sec:miss-slutsats}

Ja, med förbehåll.  Nybörjare missförstår mycket riktigt vad en klass
och ett objekt är, och på just de sätt som materialets tre
tilläggsfrågor riktar sig mot: klassen som en samling av sina objekt
snarare än en mall \parencite{Ragonis2005OOP,EckerdalThune2005,Xinogalos2015},
attributen som objektets identitet \parencite{Holland1997} och en
allmänt vag bild av vad ett objekt är
\parencite{Kaczmarczyk2010,SandersThomas2007}.  Förbehållen är att
studierna gäller Java och Smalltalk, att de inte säger hur vanliga
missuppfattningarna är hos våra studenter, och att svårigheten delvis
beror på programmens storlek och undervisningens ordning
\parencite{Wiedenbeck1999,KunkleAllen2016}.  Därför ställs frågorna i
materialet som frågor, och svaret prövas i ett motexempel --- som Holland
med flera föreslår --- i stället för att missuppfattningarna påstås om
studenterna.

\section{Fortsatt arbete}
\label{sec:miss-fortsatt}

Sökningen lämnade tre saker ogjorda.  Två databaser saknas helt: IEEE
Xplore hade nått sin dagsgräns och Semantic Scholar avvisade sin nyckel
(\cref{sec:miss-diskussion}).  Samma fem frågor bör ställas där när
tjänsterna svarar igen, och även till DBLP, som svarade ojämnt; IEEE
indexerar just den konferenslitteratur om programmeringsundervisning som
frågorna söker.  Fem källor är därtill lästa bara i sammanfattning.
Tyngst väger Ragonis och Ben-Aris artikel\autocite{Ragonis2005OOP}, som
är stängd: citatet om klassen som en samling av objekt är ärvt från
författarens egen genomgång av missuppfattningar, och är därmed
kapitlets enda led som inte kontrollerats mot en källtext.  De fyra
övriga\autocite{Xinogalos2015,SandersThomas2007,Wiedenbeck1999,%
KunkleAllen2016} bär frekvenssiffran, belägget i studenternas egna
program och de två förbehållen om programstorlek och
undervisningsordning; en läsning i fulltext skulle visa hur stora
skillnaderna var och i vilka populationer de uppmättes.  Också de nio
stödjande posterna i \cref{tab:sok-missuppfattningar} är lästa bara som
sammanfattningar, och en genomläsning skulle göra antalet till en
räkning i stället för en övre gräns.

För ett säkrare svar saknas två slags studier.  Ingen av källorna gäller
Python, och ingen gäller den population materialet vänder sig till:
nybörjare som möter klasser i ett språk där attributen syns och där
klassen själv är ett objekt.  En frekvensstudie i kursens egen
population --- de tre frågorna ställda som ett diagnostiskt prov ---
skulle visa hur vanliga missuppfattningarna är här, vilket ingen av de
citerade studierna kan svara på.  Den andra luckan gäller botemedlet:
rådet att låta studenterna pröva motexempel\autocite{Holland1997} är en
erfarenhet från undervisning, inte en mätning, och ett kontrollerat
försök som jämför motexempel med att i stället berätta om
missuppfattningen skulle visa om greppet bär.  Faller ett sådant försök
ut till motexemplets fördel kan materialet använda samma grepp för fler
missuppfattningar än den som \cref{sec:dundermetoder} tar upp; faller
det ut åt andra hållet bör frågorna i
\cref{sec:med-klasser,sec:klasser-objekt} formuleras om, eftersom de då
tar tid utan att ge något.

\section{Träffar som rör påståendet}
\label{sec:miss-traffar}

\Cref{tab:sok-missuppfattningar} listar de poster som rör påståendet
--- citerade, stödjande och kvalificerande --- med skälet till varje
beslut; de angränsande och de felträffar som uteslöts anges som antal i
tabelltexten, och antalet träffar per fråga står i
\cref{tab:sok-queries-B}.
% Author's note (verification and reproduction):
%   scholar session : classes-missuppfattningar-v2 (supersedes
%                     classes-missuppfattningar, whose classification
%                     lacked the four category tags the audit table
%                     needs; kept only in the scholar session store)
%   exports         : litteratursokning/classes-missuppfattningar-v2.csv,
%                     litteratursokning/classes-missuppfattningar-v2-full.tex
%   searched        : 2026-09-03; providers openalex wos scopus dblp
%                     (dblp: HTTP 503 on the first S1 call; the
%                     three-word lists 'object-oriented novices
%                     misconceptions', 'object misconceptions students',
%                     'object-oriented programming misconceptions review',
%                     'objects-first novices comparison', 'object-oriented
%                     misconceptions replication' gave 0/2/0/0/0 hits and
%                     were rerun as 'object-oriented misconceptions',
%                     'object misconceptions', 'programming misconceptions
%                     review', 'objects-first novices', 'misconceptions
%                     replication' = 4/9/1/0/1 in the v2 session (the
%                     v1 session holds both runs); ieee: daily quota exhausted, absent;
%                     s2: key rejected HTTP 403, absent);
%                     limit 40 per provider and query; exact strings per
%                     provider in the session export and in tab:sok-queries-B
%   classification  : scholar llm context + scholar llm classify
%                     --no-examples --no-enrich, model github_copilot/gpt-5.4;
%                     human review of every qualifies-claim row and every
%                     row with confidence < 0.7 (abstracts); one abstract-less
%                     row (10.1016/j.scico.2012.10.006) discarded by hand
%   theme tags      : scholar sessions decide --keep --doi ... -t <theme>:
%                     class-vs-template, object-concept, comprehension,
%                     teaching-order
%   hit-list table  : scholar sessions export classes-missuppfattningar-v2
%                     -f table --bearing-only --lang sv
%                     --label sok-missuppfattningar
%                     --track "missuppfattningar om klasser och objekt"
%                     --theme class-vs-template="klass som mall eller samling"
%                     --theme object-concept="objektbegreppet"
%                     --theme comprehension="programförståelse"
%                     --theme teaching-order="undervisningsordning"
%                     -o litteratursokning/classes-missuppfattningar-v2-full
%   provenance      : ltnotes.bib, keys Ragonis2005OOP Holland1997
%                     Kaczmarczyk2010 (reused from vt-prog-misconceptions,
%                     re-verified here as stated in each block)
%                     EckerdalThune2005 Xinogalos2015 SandersThomas2007
%                     Wiedenbeck1999 KunkleAllen2016 (CLAIM / FOUND-VIA /
%                     PICKED / QUOTE / VERIFIED / COUNTER / DATE)
%   known items     : Ragonis2005OOP, Holland1997, Kaczmarczyk2010 come
%                     from the misconceptions paper's protocol and were
%                     not returned by these queries (not in the table)
%   full text       : Holland1997 (ORO manuscript via Wayback),
%                     Kaczmarczyk2010 (author's open copy),
%                     EckerdalThune2005 (the user's reMarkable, 2026-09-04)
%   abstract-level  : Ragonis2005OOP (closed; quote inherited),
%                     Xinogalos2015, SandersThomas2007,
%                     KunkleAllen2016 (ACM refuses automated download),
%                     Wiedenbeck1999 (Elsevier, closed) -- for the user
%                     to verify
% --- scholar LaTeX fragment --------------------------------------------
% Generated by: scholar sessions export -f table
% Session: classes-missuppfattningar-v2
% \input this file from the body of a document whose
% preamble loads:
%   \usepackage{longtable}
%   \usepackage{booktabs}
%   \usepackage{array}
% ----------------------------------------------------------------------

\begingroup\footnotesize
\begin{longtable}{@{}%
  >{\raggedright\arraybackslash}p{0.44\textwidth}c%
  >{\raggedright\arraybackslash}p{0.13\textwidth}%
  >{\raggedright\arraybackslash}p{0.26\textwidth}@{}}
\caption{Träffar som rör påståendet, missuppfattningar om klasser och objekt (6 citerade, 9 som stöder påståendet och 4 som kvalificerar eller motsäger det, av 387 arbeten; de 152 angränsande och 216 felträffarna listas inte).  Ett arbete som flera databaser lämnat står på en rad, med databaserna uppräknade. Skälet till varje inkludering står i sista kolumnen; för maskinklassade rader anges modellens konfidens. Databaser: OpenAlex (OA; inte open access), DBLP, Web of Science (WoS), Scopus.}\label{tab:sok-missuppfattningar}\\
\toprule
Titel & År & Leverantör & Skäl \\
\midrule
\endfirsthead
\multicolumn{4}{@{}l}{\emph{\tablename~\thetable{} (forts.)}}\\
\toprule
Titel & År & Leverantör & Skäl \\
\midrule
\endhead
\midrule \multicolumn{4}{r@{}}{\emph{forts.\ på nästa sida}}\\
\endfoot
\bottomrule
\endlastfoot
A comparison of the comprehension of object-oriented and procedural programs by novice programmers & 1999 & WoS, OA & \emph{citerad}: programförståelse (oo-vs-procedural; program-comprehension) \\
Avoiding object misconceptions. & 1997 & DBLP & \emph{citerad}: attributen som identitet (object-misconceptions; teaching) \\
Checklists for grading object-oriented CS1 programs: concepts and misconceptions. & 2007 & OA, DBLP, WoS & \emph{citerad}: objektbegreppet (class-object-misconceptions; cs1) \\
Novice Java programmers' conceptions of \textquotedbl{}object\textquotedbl{} and \textquotedbl{}class\textquotedbl{}, and variation theory & 2005 & OA, Scopus & \emph{citerad}: klass som mall eller samling (object-class-conceptions; phenomenography) \\
Object-Oriented Design and Programming: An Investigation of Novices' Conceptions on Objects and Classes & 2015 & WoS & \emph{citerad}: klass som mall eller samling (object-class-conceptions; undergraduates) \\
The Impact of Different Teaching Approaches and Languages on Student Learning of Introductory Programming Concepts & 2016 & OA & \emph{citerad}: undervisningsordning (teaching-approaches; language-comparison) \\
A Diagnostic Tool for Assessing Students' Perceptions and Misconceptions Regards the Current Object \textquotedbl{}this\textquotedbl{}. & 2018 & DBLP & stöder påståendet (0.93) (this-reference; diagnostic-tool) \\
Exploring the Difficulties of Learning Object-Oriented Techniques & 1997 & OA & stöder påståendet (0.95) (basic-object-concepts; learning-difficulties) \\
Object-Oriented Programming: Diagnosis Understanding by Identifying and Describing Novice Perceptions & 2020 & WoS, Scopus & stöder påståendet (0.90) (misconceptions; concept-inventory) \\
On the Students' Misconceptions in Object-Oriented Language Constructs. & 2019 & DBLP & stöder påståendet (0.87) (language-constructs; misconceptions) \\
Statistical frequency-analysis of misconceptions in object-oriented-programming - Regularized PCR models for frequency analysis across OOP concepts and related factors & 2018 & Scopus & stöder påståendet (0.81) (misconception-frequency; oop-concepts) \\
Students' Conceptions of Object-Oriented Programming in the Context of Game Designing in Computing Education: Design of a Ph.D. Research Project to Explore Students' Conceptions in a Long-Term Study & 2021 & Scopus & stöder påståendet (0.80) (students-conceptions; game-design-context) \\
Students' Conceptions of Programming in the Context of Game Design & 2022 & Scopus & stöder påståendet (0.79) (students-conceptions; game-design-context) \\
Thinking in imperative or objects? A study on how novice programmer thinks when it comes to designing an application & 2019 & Scopus & stöder påståendet (0.72) (imperative-vs-objects-thinking; design-thinking) \\
Understanding the \textquotedbl{}this\textquotedbl{} reference in object oriented programming: Misconceptions, conceptions, and teaching recommendations. & 2021 & DBLP & stöder påståendet (0.96) (this-reference; misconceptions) \\
An empirical study of novice program comprehension in the imperative and object-oriented styles & 1997 & OA & kvalificerar eller motsäger påståendet (0.95) (oo-vs-imperative; novice-comprehension) \\
An examination of procedural and object-oriented systems analysis methods: Does prior experience help or hinder performance? & 1999 & OA, WoS & kvalificerar eller motsäger påståendet (0.84) (oo-vs-procedural; novice-analysts) \\
Comparison of Two Approaches when Teaching Object-Orientated Programming to Novices & 2011 & WoS & kvalificerar eller motsäger påståendet (0.98) (objects-first; teaching-order-comparison) \\
Is there any difference in novice comprehension of a small program written in the event-driven and object-oriented styles? & 2002 & Scopus & kvalificerar eller motsäger påståendet (0.72) (oo-vs-event-driven; novice-comprehension) \\
\end{longtable}
\endgroup


\chapter{Är inkapsling en etablerad princip?}
\label{app:inkapsling}
\chapterprecis{Författaren har ännu inte granskat resultaten i den här
  bilagan i sin helhet.}

I \cref{sec:tva-exempel} säger materialet att de dubbla understrecken
gör attributen privata, att det kallas inkapsling och att inkapsling är
en viktig princip i objektorienterad programmering.  Det sista är ett
faktapåstående om principens ställning, och det döljer två frågor: var
idén kommer ifrån, och om den är \emph{etablerad}.  Idén att låta en
modul dölja sina designbeslut för alla andra brukar tillskrivas
Parnas\autocite{Parnas1972Criteria}, men att en idé har ett ursprung
säger inte att den är allmänt omfattad.  Frågan som ska besvaras är
därför: \emph{är inkapsling en etablerad princip i objektorienterad
programmering, och varifrån kommer den?}  \Cref{sec:ink-metod}
beskriver sökningen, \cref{sec:ink-resultat} vad den fann,
\cref{sec:ink-diskussion} vad som begränsar svaret,
\cref{sec:ink-slutsats} ger det och \cref{sec:ink-fortsatt} säger vad
som återstår att göra; \cref{tab:sok-inkapsling} sist i
kapitlet listar de träffar som rör påståendet.

\section{Metod}
\label{sec:ink-metod}

Frågan söktes ämnesdrivet och tvåsidigt den 3 september 2026 med
verktyget \texttt{scholar}, mot de databaser som svarade: OpenAlex, Web
of Science och Scopus, samt DBLP, som svarade på en del av frågorna och
avvisade andra med ett serverfel.  IEEE Xplore hade nått sin dagsgräns
och Semantic Scholar avvisade sin nyckel; båda saknas.  Stödfrågorna
söker begreppet under dess båda namn --- \emph{encapsulation} och
\emph{information hiding} --- som designprincip i objektorienterad
programmering, i modulariseringsforskningen och i undervisningen;
motfrågorna söker med samma begrepp efter kritik, begränsningar och
avsaknad av empiriskt stöd, och därtill efter fältets egen invändning
mot hur principen brukar tillämpas: att läs- och skrivmetoder
(\foreignlanguage{english}{getters} och
\foreignlanguage{english}{setters}) urholkar den.
\Cref{tab:sok-queries-C} visar frågorna och antalet träffar per databas.
Parnas artikel från 1972 hämtades som känt dokument; den var redan läst
och verifierad för föreläsningen \emph{Funktioner}, där den bär SRP:s
härstamning, och kopian av dess anteckning står i källfilen.

% Author's note (verification and reproduction): see the block above the
% hit-list \input at the end of this chapter.  The query table is
% transcribed from the session's query records (scholar sessions show
% classes-inkapsling) and the per-provider hit counts recorded there.
\begin{table}[!htb]
  \centering
  \caption{Sökfrågor och antal träffar per databas, inkapsling,
    3 september 2026, verktyget \texttt{scholar}.  S = stödfråga, M =
    motfråga.  Databaser: OpenAlex (OA; inte
    \foreignlanguage{english}{open access}), DBLP, Web of Science (WoS),
    Scopus.  Högst 40 träffar hämtades per databas och fråga, så 40
    betyder \enquote{minst 40}.  Frågan visas i den booleska form som
    skickades till OpenAlex; Web of Science fick samma sträng som
    \texttt{TS=(\dots)}, Scopus som \texttt{TITLE-ABS-KEY(\dots)} och
    DBLP en kort ordlista med samma begrepp (DBLP indexerar bara titlar
    och kräver att varje ord förekommer; ordlistorna med tre ord gav få
    träffar, så frågorna ställdes om med två ord --- ordlistorna står i
    författarens anteckning).  IEEE Xplore hade nått sin
    dagsgräns och Semantic Scholar avvisade sin nyckel; de
    saknas.}
  \label{tab:sok-queries-C}
  \footnotesize
  \begin{tabular}{@{}lp{0.55\linewidth}rrrr@{}}
    \toprule
    & Nyckelord & OA & DBLP & WoS & Scopus \\
    \midrule
    S1 & (encapsulation \textsc{or} \enquote{information hiding}) \textsc{and} (\enquote{object-oriented} \textsc{or} \enquote{object oriented}) \textsc{and} (principle \textsc{or} principles \textsc{or} \enquote{design principle})
       & 40 & 2 & 40 & 40 \\
    S2 & (\enquote{information hiding} \textsc{or} encapsulation) \textsc{and} (modularity \textsc{or} modules \textsc{or} \enquote{modular design}) \textsc{and} (criteria \textsc{or} benefits \textsc{or} maintainability \textsc{or} \enquote{software design})
       & 40 & 5 & 40 & 40 \\
    S3 & (encapsulation \textsc{or} \enquote{information hiding} \textsc{or} \enquote{data hiding}) \textsc{and} (teaching \textsc{or} education \textsc{or} novices \textsc{or} \enquote{introductory programming}) \textsc{and} (\enquote{object-oriented} \textsc{or} classes)
       & 40 & 2 & 40 & 40 \\
    \midrule
    M1 & (encapsulation \textsc{or} \enquote{information hiding}) \textsc{and} (critique \textsc{or} criticism \textsc{or} limitations \textsc{or} harmful \textsc{or} \enquote{not beneficial} \textsc{or} myth \textsc{or} \enquote{no evidence} \textsc{or} \enquote{empirical evidence})
       & 40 & 0 & 40 & 40 \\
    M2 & (getters \textsc{or} setters \textsc{or} accessors \textsc{or} \enquote{accessor methods}) \textsc{and} (encapsulation \textsc{or} \enquote{information hiding} \textsc{or} \enquote{object-oriented}) \textsc{and} (harmful \textsc{or} evil \textsc{or} critique \textsc{or} \enquote{anti-pattern} \textsc{or} \enquote{code smell} \textsc{or} violate)
       & 40 & 1 & 0 & 1 \\
    \bottomrule
  \end{tabular}
\end{table}

Sammanlagt lämnade databaserna 448 poster, vilket efter samma
sammanslagning av dubbletter som i \cref{app:missuppfattningar} blir 435
arbeten.  Alla klassades på samma
sätt som i \cref{app:missuppfattningar}: först med en språkmodell mot en
beskrivning av frågan och fyra kategorier, sedan för hand för varje post
som modellen satt i den kvalificerande kategorin och varje post med
konfidens under 0,7.  En rad i \cref{tab:sok-inkapsling} är rättad för
hand: OpenAlex lämnade Parnas artikel som den tekniska rapporten från
1971, medan kapitlet citerar tidskriftsversionen från 1972, och det är
den årtalet nu anger.  Ordet \emph{encapsulation} betyder annat i andra
fält --- solceller, läkemedel, nätverkspaket --- och drygt hundra poster
föll bort av det skälet.  De källor som kapitlet citerar lästes i
fulltext där den gick att nå; för de övriga står det i källfilen att
bara sammanfattningen lästs.

\section{Resultat}
\label{sec:ink-resultat}

\paragraph{Ursprung.}  Idén härrör från Parnas, som beskriver den så:
\textcquote[1056]{Parnas1972Criteria}{Every module in the second
decomposition is characterized by its knowledge of a design decision
which it hides from all others}.  Att det är dit fältet spårar idén
bekräftar en senare forskningsartikel i förbigående:
\textcquote{Sullivan2001Modularity}{In 1972 Parnas introduced information
hiding as an approach to devising modular structures for software
designs}.  Ordet \emph{encapsulation} för samma idé i objektorienterade
språk definierar \textcite{Snyder1986}:
\textcquote{Snyder1986}{Encapsulation is a technique for minimizing
interdependencies among separately-written modules by defining strict
external interfaces}, och \textcquote{Snyder1986}{[a] module is
encapsulated if clients are restricted by the definition of the
programming language to access the module only via its defined external
interface} --- det är det de dubbla understrecken gör i
\cref{sec:tva-exempel}.

\paragraph{Hur etablerad?}  Mycket.  \Textcite{Sullivan2001Modularity}
inleder med att \textcquote{Sullivan2001Modularity}{[t]he concept of
information hiding modularity is a cornerstone of modern software design
thought}; \textcite{Tempero2009Fields} med att
\textcquote{Tempero2009Fields}{[t]he information hiding principle is
generally accepted as one that i[f] followed leads to higher quality
software than if it is not followed}; och kritikerna
\textcite{Ostermann2011Revisiting} med att informationsdöljande i
forskarsamhället \textcquote{Ostermann2011Revisiting}{is nowadays such
an undisputed dogma of modularity that Fred Brooks even felt that he had
to apologize to Parnas for questioning it}.  Sökningen ger därtill ett
femtontal poster --- läroböcker, mätningar av designkvalitet, formella
grunder --- som tar principen för given; de står som stödjande i
\cref{tab:sok-inkapsling}.

\paragraph{Kritik och förbehåll.}  Motsökningen bär tre förbehåll.  Det
första är gammalt: \textcite{Snyder1986} visar i samma artikel att
\textcquote{Snyder1986}{in most object-oriented languages, the
introduction of inheritance severely compromises encapsulation}, eftersom
en underklass som får läsa förälderns attribut gör dem till en del av
kontraktet --- vilket är skälet till att \mintinline{python}{PersonNick}
i \cref{sec:arv} går via \mintinline{python}{get_first_name()} och
\mintinline{python}{super()} i stället för att röra förälderns attribut.
Det andra är teoretiskt: samma \textcite{Sullivan2001Modularity} som
kallar principen en hörnsten säger att
\textcquote{Sullivan2001Modularity}{its formulation remains casual}, och
\textcite{Ostermann2011Revisiting} menar att dogmen bör omprövas ---
dock \textcquote{Ostermann2011Revisiting}{[w]e do of course not propose
to abandon information hiding or modularity in general}.  Det tredje är
praktiskt: principen följs ofullständigt.  \Textcite{Tempero2009Fields}
fann i hundra öppna Javaprogram att
\textcquote{Tempero2009Fields}{it is not uncommon (albeit not that
terribly common) to declare non-private fields}, och
\textcite{ZollerSchmolitzky2012} att
\textcquote{ZollerSchmolitzky2012}{around one third of all type and
method access modifiers} är generösare än nödvändigt.  Invändningen att
läs- och skrivmetoder urholkar inkapslingen --- som materialet självt
använder --- finns knappt i den indexerade litteraturen: motfrågan M2 gav
noll träffar i Web of Science och en i Scopus, och den kända kritiken är
inte vetenskapligt publicerad, så den citeras inte.

\section{Diskussion och begränsningar}
\label{sec:ink-diskussion}

Svaret vilar på fyra källor lästa i fulltext (Parnas, Snyder, Sullivan
med flera, Ostermann med flera) och två lästa i sammanfattning (Tempero;
Zoller och Schmolitzky), som är stängda hos IEEE; de senare bär
förbehållet om praktiken och bör kontrolleras av den som har tillgång.
Frågan är ovanligt bred --- \emph{encapsulation} och \emph{information
hiding} förekommer i tusentals titlar --- så de fyrtio träffar som
hämtades per fråga och databas är ett urval, inte en totalräkning, och
antalet stödjande poster i \cref{tab:sok-inkapsling} är därför en undre
gräns för stödet och ingen mätning av det.  IEEE Xplore och Semantic
Scholar saknas, DBLP svarade på få ord i taget, och klassningen är
modellstödd.  Slutligen gäller de empiriska källorna Java, inte Python,
där inkapslingen dessutom är en konvention snarare än ett språkkrav:
ett attribut med dubbla understreck skrivs om i stället för att skyddas,
som fotnoten om namnmangling i \cref{sec:tva-exempel} beskriver.

\section{Slutsats}
\label{sec:ink-slutsats}

Ja, med förbehåll.  Inkapsling är en etablerad princip i objektorienterad
programmering --- \enquote{a cornerstone}, \enquote{generally accepted},
till och med \enquote{an undisputed dogma}
\parencite{Sullivan2001Modularity,Tempero2009Fields,Ostermann2011Revisiting}
--- och idén om att dölja designbeslut härrör från
\textcite{Parnas1972Criteria}.  Förbehållen är att principen är löst
formulerad och i sin klassiska form omdiskuterad, att arv kan bryta den,
och att den i praktiken följs ofullständigt.  Materialet använder den i
sin enklaste form --- privata attribut och läsmetoder --- och säger
\enquote{viktig princip}, inte \enquote{alltid rätt}.

\section{Fortsatt arbete}
\label{sec:ink-fortsatt}

Sökningen lämnade tre saker ogjorda.  IEEE Xplore hade nått sin
dagsgräns och Semantic Scholar avvisade sin nyckel, och för den här
frågan är det en dubbel lucka: de två källor som bär förbehållet om
praktiken\autocite{Tempero2009Fields,ZollerSchmolitzky2012} är stängda
just hos IEEE och lästa bara i sammanfattning
(\cref{sec:ink-diskussion}).  Fulltexterna skulle visa vad de räknade
som ett fält respektive en åtkomstmodifierare, i vilka program och med
vilket urval, och därmed hur långt siffrorna räcker.  Motfrågan om
läs- och skrivmetoder --- M2 i \cref{tab:sok-queries-C} --- gav noll
träffar i Web of Science och en i Scopus: invändningen finns, men inte i
den indexerade tidskriftslitteraturen, och den som vill pröva den får
söka i böcker och fackpress i stället.  Slutligen är frågan så bred att
de fyrtio träffarna per fråga och databas är ett urval; en körning utan
tak skulle ge stödets storlek och inte bara dess undre gräns.

För ett säkrare svar saknas mätningar på nybörjare och på Python.  De
källor som visar att principen är etablerad är designlitteratur, och de
två empiriska gäller professionellt skriven Javakod; ingen av dem prövar
vad inkapslingen ger den som just har lärt sig skriva klasser.
Delpåståendet som ingen har prövat är alltså det materialet självt
använder: att ett program med privata attribut och läsmetoder är
lättare att ändra än ett med öppna attribut.  Ett kontrollerat försök
där nybörjare ändrar samma Pythonprogram i två versioner, en med och en
utan privata attribut, skulle avgöra saken.  Visar försöket en vinst kan
materialet säga att principen är prövad och inte bara etablerad; visar
det ingen vinst för små program bör texten säga att principen betalar
sig först när programmen växer.

\section{Träffar som rör påståendet}
\label{sec:ink-traffar}

\Cref{tab:sok-inkapsling} listar de poster som rör påståendet ---
citerade, stödjande och kvalificerande --- med skälet till varje beslut;
de angränsande och de felträffar som uteslöts anges som antal i
tabelltexten, och antalet träffar per fråga står i
\cref{tab:sok-queries-C}.
% Author's note (verification and reproduction):
%   scholar session : classes-inkapsling-v2 (supersedes
%                     classes-inkapsling, whose classification lacked
%                     the four category tags the audit table needs; kept
%                     only in the scholar session store)
%   exports         : litteratursokning/classes-inkapsling-v2.csv,
%                     litteratursokning/classes-inkapsling-v2-full.tex
%   searched        : 2026-09-03; providers openalex wos scopus dblp
%                     (dblp: the three-word lists 'encapsulation
%                     object-oriented principle', 'information hiding
%                     modularity', 'encapsulation teaching
%                     object-oriented', 'encapsulation criticism',
%                     'getters setters encapsulation' gave 0/5/0/0/0 hits
%                     and were rerun as 'encapsulation principle',
%                     'information hiding modularity', 'encapsulation
%                     teaching', 'information hiding critique', 'getters
%                     setters' = 2/5/2/0/1 in the v2 session (the v1
%                     session holds both runs); ieee: daily quota exhausted, absent; s2: key
%                     rejected HTTP 403, absent); limit 40 per provider
%                     and query; exact strings per provider in the
%                     session export and in tab:sok-queries-C
%   classification  : scholar llm context + scholar llm classify
%                     --no-examples --no-enrich, model github_copilot/gpt-5.4;
%                     human review of every qualifies-claim row and every
%                     row with confidence < 0.7 (abstracts)
%   theme tags      : scholar sessions decide --keep --doi ... -t <theme>:
%                     origin, definition, adoption, practice, critique
%   hit-list table  : scholar sessions export classes-inkapsling-v2 -f table
%                     --bearing-only --lang sv --label sok-inkapsling
%                     --track "inkapsling" --theme origin="ursprung"
%                     --theme definition="definition"
%                     --theme adoption="etablering" --theme practice="praktik"
%                     --theme critique="kritik"
%                     -o litteratursokning/classes-inkapsling-v2-full
%   provenance      : ltnotes.bib, keys Parnas1972Criteria (reused from
%                     Funktioner) Snyder1986 Sullivan2001Modularity
%                     Tempero2009Fields ZollerSchmolitzky2012
%                     Ostermann2011Revisiting (CLAIM / FOUND-VIA / PICKED /
%                     QUOTE / VERIFIED / COUNTER / DATE)
%   known items     : Parnas1972Criteria (surfaced as its 1971 technical-
%                     report record, doi 10.21236/ad0773837, which is the
%                     row tagged "ursprung" in the table)
%   full text       : Snyder1986 (Tufts course copy), Sullivan2001Modularity
%                     (Drexel co-author copy), Ostermann2011Revisiting
%                     (CMU co-author copy)
%   abstract-level  : Tempero2009Fields, ZollerSchmolitzky2012 (IEEE,
%                     closed) -- for the user to verify
% --- scholar LaTeX fragment --------------------------------------------
% Generated by: scholar sessions export -f table
% Session: classes-inkapsling-v2
% \input this file from the body of a document whose
% preamble loads:
%   \usepackage{longtable}
%   \usepackage{booktabs}
%   \usepackage{array}
% ----------------------------------------------------------------------

\begingroup\footnotesize
\begin{longtable}{@{}%
  >{\raggedright\arraybackslash}p{0.44\textwidth}c%
  >{\raggedright\arraybackslash}p{0.13\textwidth}%
  >{\raggedright\arraybackslash}p{0.26\textwidth}@{}}
\caption{Träffar som rör påståendet, inkapsling (6 citerade, 15 som stöder påståendet och 0 som kvalificerar eller motsäger det, av 439 unika träffar; de 170 angränsande och 248 felträffarna listas inte). Skälet till varje inkludering står i sista kolumnen; för maskinklassade rader anges modellens konfidens. Databaser: OpenAlex (OA; inte open access), DBLP, Web of Science (WoS), Scopus.}\label{tab:sok-inkapsling}\\
\toprule
Titel & År & Leverantör & Skäl \\
\midrule
\endfirsthead
\multicolumn{4}{@{}l}{\emph{\tablename~\thetable{} (forts.)}}\\
\toprule
Titel & År & Leverantör & Skäl \\
\midrule
\endhead
\midrule \multicolumn{4}{r@{}}{\emph{forts.\ på nästa sida}}\\
\endfoot
\bottomrule
\endlastfoot
Encapsulation and inheritance in object-oriented programming languages & 1986 & OA & \emph{citerad}: definition (inheritance-vs-encapsulation; requirements-for-support) \\
How Fields are Used in Java: An Empirical Study & 2009 & WoS & \emph{citerad}: praktik (empirical-practice; non-private-fields) \\
Measuring Inappropriate Generosity with Access Modifiers in Java Systems & 2012 & WoS & \emph{citerad}: praktik (empirical-practice; access-modifiers; too-generous-access) \\
On the Criteria to Be Used in Decomposing Systems into Modules & 1971 & OA & \emph{citerad}: ursprung (parnas; information-hiding) \\
Revisiting Information Hiding: Reflections on Classical and Nonclassical Modularity. & 2011 & DBLP & \emph{citerad}: kritik (revisiting-information-hiding; classical-vs-nonclassical-modularity) \\
The structure and value of modularity in software design & 2001 & OA & \emph{citerad}: etablering (information-hiding; parnas; modularity-benefits) \\
A new approach to component-oriented programming & 2004 & WoS & stöder påståendet (0.94) (component-programming; encapsulation-key-principle) \\
Architectural design methodologies for complex evolving systems & 2007 & WoS & stöder påståendet (0.91) (encapsulation-benefits; evolvability) \\
Assuring good style for object-oriented programs & 1989 & OA & stöder påståendet (0.99) (law-of-demeter; information-hiding; modularity) \\
Concepts and paradigms of object-oriented programming & 1990 & OA & stöder påståendet (0.92) (origins; oop-paradigms; robust-paradigm) \\
Data abstraction and information hiding & 2002 & OA & stöder påståendet (0.98) (data-abstraction; information-hiding; modular-soundness) \\
Designing information hiding modularity for model transformation languages. & 2014 & DBLP & stöder påståendet (0.94) (information-hiding; modularity; model-transformation) \\
Encapsulation, Reusability and Extensibility in Object-Oriented Programming Languages & 1987 & OA & stöder påståendet (0.99) (comparative-survey; strict-information-hiding) \\
Fundamentals of object-oriented design in UML & 1999 & OA & stöder påståendet (0.98) (textbook-like; history; parnas) \\
Information hiding interfaces for aspect-oriented design & 2005 & OA & stöder påståendet (0.97) (information-hiding; aop-interfaces; modularity) \\
Interfaces for strongly-typed object-oriented programming & 1989 & OA & stöder påståendet (0.93) (interfaces; design-maintenance) \\
Object-Oriented Languages & 1997 & OA & stöder påståendet (0.99) (information-hiding; oop-support) \\
Reusing code for modernization of legacy systems & 2006 & WoS & stöder påståendet (0.92) (information-hiding; legacy-modernization) \\
Simple type-theoretic foundations for object-oriented programming & 1994 & OA & stöder påståendet (0.98) (formal-foundations; definition-of-encapsulation) \\
The Modular Structure of Complex Systems & 1985 & OA & stöder påståendet (1.00) (parnas; information-hiding; module-guide) \\
The Opaque Pointer Design Pattern in Python: Towards a Pythonic PIMPL for Modularity, Encapsulation, and Stability & 2026 & Scopus & stöder påståendet (0.86) (design-pattern; modularity; encapsulation) \\
\end{longtable}
\endgroup


\chapter{Bör man föredra komposition framför arv?}
\label{app:komposition}
\chapterprecis{Författaren har ännu inte granskat resultaten i den här
  bilagan i sin helhet.}

I \cref{sec:arv} ger materialet tumregeln att föredra komposition
framför arv och att använda arv bara vid en verklig
\foreignlanguage{english}{is-a}-relation, och tillskriver rådet Gamma,
Helm, Johnson och Vlissides\autocite{Gamma1994DesignPatterns}.
Det är tre påståenden: att rådet kommer därifrån, att det är etablerat i
objektorienterad design, och --- underförstått --- att det är riktigt,
det vill säga att arv har kostnader som komposition slipper.  Frågan
som ska besvaras är: \emph{bör man föredra komposition framför arv, och
kommer rådet från Gamma med flera?}  \Cref{sec:komp-metod} beskriver
sökningen, \cref{sec:komp-resultat} vad den fann,
\cref{sec:komp-diskussion} vad som begränsar svaret,
\cref{sec:komp-slutsats} ger det och \cref{sec:komp-fortsatt} säger vad
som återstår att göra; \cref{tab:sok-komposition} sist i
kapitlet listar de träffar som rör påståendet.

\section{Metod}
\label{sec:komp-metod}

Frågan söktes ämnesdrivet och tvåsidigt den 3 september 2026 med
verktyget \texttt{scholar}, mot de databaser som svarade: OpenAlex, Web
of Science, Scopus och DBLP.  IEEE Xplore hade nått sin dagsgräns och
Semantic Scholar avvisade sin nyckel; båda saknas.  Stödfrågorna söker
rådet under dess egna namn (\emph{composition over inheritance},
\emph{favor composition}), dess mekanismer (\emph{object composition},
\emph{delegation}) som designprincip, och arvets belagda problem
(\emph{fragile base class}, koppling, missbruk); motfrågorna söker med
samma begrepp efter arvets fördelar och empiriska värde, och efter
kritik av rådet och av delegering.  \Cref{tab:sok-queries-D} visar
frågorna och antalet träffar per databas.  Boken \emph{Design
Patterns} är inte indexerad i databaserna och finns inte öppet
tillgänglig; den lästes i stället som känt dokument i författarens eget
exemplar, och sidan och den ordagranna lydelsen står i källfilen.

% Author's note (verification and reproduction): see the block above the
% hit-list \input at the end of this chapter.  The query table is
% transcribed from the session's query records (scholar sessions show
% classes-komposition) and the per-provider hit counts recorded there.
\begin{table}[!htb]
  \centering
  \caption{Sökfrågor och antal träffar per databas,
    komposition framför arv, 3 september 2026, verktyget
    \texttt{scholar}.  S = stödfråga, M = motfråga.  Databaser: OpenAlex
    (OA; inte \foreignlanguage{english}{open access}), DBLP, Web of
    Science (WoS), Scopus.  Högst 40 träffar hämtades per databas och
    fråga, så 40 betyder \enquote{minst 40}.  Frågan visas i den
    booleska form som skickades till OpenAlex; Web of Science fick samma
    sträng som \texttt{TS=(\dots)}, Scopus som
    \texttt{TITLE-ABS-KEY(\dots)} och DBLP en kort ordlista med samma
    begrepp (DBLP indexerar bara titlar och kräver att varje ord
    förekommer; ordlistorna står i författarens anteckning).  IEEE
    Xplore hade nått sin dagsgräns och Semantic Scholar avvisade sin
    nyckel; de saknas.}
  \label{tab:sok-queries-D}
  \footnotesize
  \begin{tabular}{@{}lp{0.55\linewidth}rrrr@{}}
    \toprule
    & Nyckelord & OA & DBLP & WoS & Scopus \\
    \midrule
    S1 & (\enquote{composition over inheritance} \textsc{or} \enquote{favor composition} \textsc{or} \enquote{favour composition} \textsc{or} \enquote{prefer composition}) \textsc{and} (inheritance \textsc{or} \enquote{object-oriented} \textsc{or} \enquote{object oriented})
       & 40 & 17 & 2 & 3 \\
    S2 & (inheritance \textsc{or} \enquote{class inheritance}) \textsc{and} (\enquote{object composition} \textsc{or} delegation) \textsc{and} (\enquote{design principle} \textsc{or} \enquote{design patterns} \textsc{or} reuse \textsc{or} \enquote{software design})
       & 40 & 0 & 30 & 40 \\
    S3 & inheritance \textsc{and} (\enquote{object-oriented} \textsc{or} \enquote{object oriented}) \textsc{and} (problems \textsc{or} harmful \textsc{or} \enquote{fragile base class} \textsc{or} coupling \textsc{or} misuse \textsc{or} \enquote{considered harmful})
       & 40 & 0 & 40 & 40 \\
    \midrule
    M1 & inheritance \textsc{and} (\enquote{object-oriented} \textsc{or} \enquote{object oriented}) \textsc{and} (benefits \textsc{or} beneficial \textsc{or} advantages \textsc{or} \enquote{code reuse} \textsc{or} empirical) \textsc{and} (maintainability \textsc{or} quality \textsc{or} defects \textsc{or} \enquote{no evidence})
       & 40 & 6 & 40 & 40 \\
    M2 & (\enquote{composition over inheritance} \textsc{or} \enquote{favor composition} \textsc{or} \enquote{favour composition} \textsc{or} delegation) \textsc{and} (critique \textsc{or} criticism \textsc{or} limitations \textsc{or} overhead \textsc{or} \enquote{not always} \textsc{or} misuse \textsc{or} drawbacks)
       & 40 & 4 & 40 & 40 \\
    \bottomrule
  \end{tabular}
\end{table}

Sammanlagt lämnade databaserna 417 poster.  Flera databaser lämnar
samma arbete, och två poster räknas som ett arbete när titlarna är
desamma bortsett från versaler och skiljetecken; efter den
sammanslagningen återstår 394 arbeten, och det är arbeten som räknas här
och i \cref{tab:sok-komposition}.  Alla klassades: först med
en språkmodell mot en beskrivning av frågan och fyra kategorier ---
stöder påståendet, kvalificerar eller motsäger det, angränsande delämne,
annat ämne --- och sedan för hand: varje post som modellen satt i den
andra kategorin, och varje post med konfidens under 0,7, lästes i
sammanfattning.  Orden \emph{inheritance} och \emph{composition} betyder
annat i genetik, juridik, musik och kemi, och ett hundratal poster föll
bort av det skälet.  De källor som kapitlet citerar kunde bara läsas i
sammanfattning, med ett undantag; det står i källfilen, och
\cref{sec:komp-diskussion} återkommer till det.

\Cref{tab:sok-komposition} rättades för hand vid genomläsningen, och
rättelserna redovisas därför att en rättelse man inte kan kontrollera
inte är någon rättelse.  Bettini med fleras artikel om dynamiska
objektbeteenden hade fått en fotnot om forskningsanslag inklistrad efter
titeln, och den avhuggen av exportens teckengräns; titeln står nu ensam.
Ett arbete om ändliga element stod på tre rader, en per databaspost, och
står nu på en.  Och en rad ströks: en \foreignlanguage{english}{narrative
review} av entitet-komponent-system i spelmotorer, utgiven som
Zenodo-preprint med en författarlista som inte går att tyda, hade
klassats som stödjande av modellen men vid genomgången bedömts som
felträff --- den var den enda rad i tabellen som saknade konfidenssiffra,
och det var därför.

\section{Resultat}
\label{sec:komp-resultat}

\paragraph{Ursprung.}  Boken \emph{Design
Patterns}\autocite{Gamma1994DesignPatterns} är den källa både
uppslagsverk och forskningsartiklar i träfflistan pekar på för rådet,
och den säger det själv, som sin \enquote{second principle of
object-oriented design}: \textcquote[20]{Gamma1994DesignPatterns}{Favor
object composition over class inheritance}.  Ingen annan upphovsman till
rådet kom fram i sökningen.  Boken är samtidigt mindre kategorisk än
rådet brukar citeras: stycket efter principen säger att
\textcquote[20]{Gamma1994DesignPatterns}{[i]nheritance and object
composition thus work together}, och att skälet till rådet är
erfarenhet --- \textcquote[20]{Gamma1994DesignPatterns}{our experience
is that designers overuse inheritance as a reuse technique}.

\paragraph{Arvets kostnader.}  Rådet vilar på att arv har kostnader
som komposition slipper, och de är belagda från flera håll.
\Textcite{Snyder1986} visade redan innan boken att
\textcquote{Snyder1986}{in most object-oriented languages, the
introduction of inheritance severely compromises encapsulation}:
en underklass som får läsa förälderns attribut gör dem till en del av
förälderns kontrakt, och föräldern kan inte längre ändras fritt.
\Textcite{WildeHuitt1992} fann att arv och polymorfism
\textcquote{WildeHuitt1992}{introduce difficulties in program analysis
and understanding} vid underhåll, och \textcite{CartwrightShepperd2000}
mätte i ett industriellt system att klasser i arvsstrukturer var
\textcquote{CartwrightShepperd2000}{approximately three times more
defect-prone} än andra.  Sökningen ger därtill ett tjugotal poster som
tar rådet för givet eller bygger vidare på det --- designmönster som
Bridge och Decorator, omstrukturering av arv till komposition --- och
de står som stödjande i \cref{tab:sok-komposition}.

\paragraph{Motsökningen.}  Här bär motsökningen kapitlet, i tre
förbehåll.  De kontrollerade experimenten är oeniga: \textcite{Daly1996}
fann att försökspersoner underhöll program med tre arvsnivåer
\textcquote{Daly1996}{significantly quicker} än en platt motsvarighet,
medan fem nivåer tog längre tid; \textcite{Harrison2000} fann tvärtom
att \textcquote{Harrison2000}{the systems without inheritance were
easier to modify than the corresponding systems containing three or
five levels of inheritance}, och efterlyste mer forskning om arvets
fördelar.  Den lärobokskostnad som brukar anföras, den sköra basklassen
(\foreignlanguage{english}{fragile base class}), visade sig i en stor
empirisk prövning
\textcquote{Sabane2017}{neither more likely to change nor more likely
to have faults}, och ingen av 104 analyserade programfel hade med den att
göra.  Och arv används mycket i verklig kod: \textcite{StevensonWood2018}
fann i 2\,440 arvshierarkier att
\textcquote{StevensonWood2018}{inheritance is very widely used but that
most of the usage patterns that occur in practice are simple in
structure}.  Arton ytterligare arbeten kvalificerar på liknande sätt
--- arvets energieffektivitet, alternativ till arv för flerfaldigt arv,
arv i testkod --- och står som kvalificerande i
\cref{tab:sok-komposition}.

\section{Diskussion och begränsningar}
\label{sec:komp-diskussion}

Flera saker begränsar svaret.  Verifieringen är ojämn: boken som rådet
tillskrivs och Snyders artikel lästes i fulltext, men de sex empiriska
studierna är stängda hos IEEE, Elsevier, Springer och ACM och lästes i
sammanfattning; det står i källfilen och bör kontrolleras av den som har
tillgång.  Källorna gäller C++, Java och
Smalltalk, inte Python, och de mäter olika saker --- underhållstid i
laboratorium, feltäthet i fält, förekomst av en viss struktur --- så
att de inte kan läggas ihop till ett mått.  Och täckningen är
ofullständig: IEEE Xplore och Semantic Scholar saknas, DBLP svarade på
få ord i taget, träffantalen i \cref{tab:sok-queries-D} är tak, och
klassningen är modellstödd.

\section{Slutsats}
\label{sec:komp-slutsats}

Ja på båda frågorna, med förbehåll.  Rådet att föredra komposition framför arv
kommer från \emph{Design Patterns}, där det står som bokens andra designprincip
(s.~20), och är etablerat i objektorienterad design: ett tjugotal poster bygger
på det, och arvets kostnader --- bruten inkapsling, svårare förståelse, fler
fel --- är belagda
\parencite{Snyder1986,WildeHuitt1992,CartwrightShepperd2000}.  Förbehållen kom
ur motsökningen: experimenten om arv och underhåll pekar åt olika håll
\parencite{Daly1996,Harrison2000}, den sköra basklassen ställer sällan till det
i praktiken \parencite{Sabane2017}, och arv används flitigt och för det mesta
enkelt \parencite{StevensonWood2018}.  Därför säger materialets tumregel
\enquote{föredra komposition, och använd arv när du verkligen har en
\foreignlanguage{english}{is-a}-relation} --- inte \enquote{undvik arv}.

\section{Fortsatt arbete}
\label{sec:komp-fortsatt}

Sökningen lämnade två databaser obesökta: IEEE Xplore hade nått sin
dagsgräns och Semantic Scholar avvisade sin nyckel, och de fem frågorna
bör ställas där när tjänsterna svarar igen.  Tyngre väger att hela den
empiriska grunden är läst i sammanfattning
(\cref{sec:komp-diskussion}): de sex studier som bär både arvets
kostnader och motsökningens förbehåll\autocite{WildeHuitt1992,%
CartwrightShepperd2000,Daly1996,Harrison2000,Sabane2017,%
StevensonWood2018} är stängda hos IEEE, Elsevier, Springer och ACM.
Två av dem avgör kapitlets viktigaste förbehåll: Daly med flera och
Harrison med flera kommer till motsatta slutsatser om arv och
underhållbarhet, och först fulltexterna säger om de mätte samma sak ---
samma uppgift, samma slags försökspersoner, samma mått --- eller om
oenigheten följer av att försöken var olika upplagda.  Att läsa dem är
det enskilt mest värdefulla som återstår.

Rådets eget namn är dessutom knappt indexerat: stödfrågan S1 gav två
träffar i Web of Science och tre i Scopus, medan samma fråga gav fyrtio
i OpenAlex och sjutton i DBLP (\cref{tab:sok-queries-D}).  Rådet lever i
böcker och i praktikerlitteratur, så den som vill mäta hur etablerat det
är får söka där.  Ursprungsfrågan är däremot färdig: boken lästes som
känt dokument, och sidan och den ordagranna lydelsen står i källfilen.

För ett säkrare svar saknas den jämförelse som rådet självt gör.  De två
oeniga experimenten varierar arvets djup mot ett platt program; ingen av
de studier som kom fram ställer samma uppgift löst med arv mot samma
uppgift löst med komposition och mäter vad det kostar att ändra den.
Det är rådets egen fråga, och den är obesvarad.  Ingen av källorna
gäller heller Python eller nybörjare: de mäter professionellt skriven
kod i C++, Java och Smalltalk.  Ett kontrollerat försök där
försökspersoner ändrar två versioner av samma program, en med arv och en
med komposition, tillsammans med en replikation av de två oeniga
experimenten under samma uppläggning, skulle avgöra saken.  Utfaller det
till kompositionens fördel kan materialet säga att tumregeln är prövad
och inte bara grundad på erfarenhet; utfaller det åt andra hållet måste
\cref{sec:arv} säga rent ut att rådet vilar på bokens
författares erfarenhet och på arvets kostnader, inte på en uppmätt
fördel för komposition.

\section{Träffar som rör påståendet}
\label{sec:komp-traffar}

\Cref{tab:sok-komposition} listar de poster som rör påståendet ---
citerade, stödjande och kvalificerande --- med skälet till varje beslut;
de angränsande och de felträffar som uteslöts anges som antal i
tabelltexten, och antalet träffar per fråga står i
\cref{tab:sok-queries-D}.
% Author's note (verification and reproduction):
%   moved           : 2026-09-22 from modules/classes/slides-even-more
%                     (Praktiska tillämpningar av klasser, bilaga B there),
%                     with its exports and bib entries, because this deck
%                     is the first in the reading order that makes the
%                     claim.  Session names and export paths unchanged;
%                     the query table's label is tab:sok-queries-D here.
%   scholar session : classes-komposition (plus classes-komposition-kanda,
%                     a throwaway known-item session used only to read
%                     four closed papers' abstracts from Web of Science)
%   exports         : litteratursokning/classes-komposition.csv,
%                     litteratursokning/classes-komposition-full.tex
%   searched        : 2026-09-03; providers openalex wos scopus dblp
%                     (dblp word lists: 'composition inheritance',
%                     'inheritance delegation reuse', 'inheritance
%                     harmful', 'inheritance maintainability',
%                     'composition critique' = 17/0/0/6/4 hits; ieee:
%                     daily quota exhausted, absent; s2: key rejected
%                     HTTP 403, absent); limit 40 per provider and
%                     query; exact strings per provider in the session
%                     export and in tab:sok-queries-D
%   classification  : scholar llm context + scholar llm classify
%                     --no-examples --no-enrich, model github_copilot/gpt-5.4;
%                     human review of every qualifies-claim row and every
%                     row with confidence < 0.7 (abstracts); two Zenodo
%                     self-published rows with bogus author fields
%                     discarded by hand as false hits
%   theme tags      : scholar sessions decide --keep --doi ... -t <theme>:
%                     costs, empirics, practice
%   hit-list table  : scholar sessions export classes-komposition -f table
%                     --bearing-only --lang sv --label sok-komposition
%                     --track "komposition framför arv"
%                     --theme costs="arvets kostnader" --theme empirics="empiri"
%                     --theme practice="praktik"
%                     -o litteratursokning/classes-komposition-full
%   provenance      : ltnotes.bib, keys Gamma1994DesignPatterns Snyder1986
%                     (this deck's bilaga C) WildeHuitt1992
%                     CartwrightShepperd2000 Daly1996 Harrison2000
%                     Sabane2017 StevensonWood2018 (CLAIM / FOUND-VIA /
%                     PICKED / QUOTE / VERIFIED / COUNTER / DATE)
%   known items     : Gamma1994DesignPatterns (a book; not indexed; read
%                     in the author's own copy on the reMarkable,
%                     2026-09-04 -- p. 20 quoted verbatim, see ltnotes.bib)
%   full text       : Snyder1986 (Tufts course copy)
%   abstract-level  : WildeHuitt1992, CartwrightShepperd2000 (IEEE),
%                     Daly1996, Sabane2017 (Springer), Harrison2000
%                     (Elsevier), StevensonWood2018 (ACM) -- for the user
%                     to verify
% --- scholar LaTeX fragment --------------------------------------------
% Generated by: scholar sessions export -f table
% Session: classes-komposition
% \input this file from the body of a document whose
% preamble loads:
%   \usepackage{longtable}
%   \usepackage{booktabs}
%   \usepackage{array}
% ----------------------------------------------------------------------

\begingroup\footnotesize
\begin{longtable}{@{}%
  >{\raggedright\arraybackslash}p{0.44\textwidth}c%
  >{\raggedright\arraybackslash}p{0.13\textwidth}%
  >{\raggedright\arraybackslash}p{0.26\textwidth}@{}}
\caption{Träffar som rör påståendet, komposition framför arv (7 citerade, 18 som stöder påståendet och 18 som kvalificerar eller motsäger det, av 394 arbeten; de 230 angränsande och 121 felträffarna listas inte).  Ett arbete som flera databaser lämnat står på en rad, med databaserna uppräknade. Skälet till varje inkludering står i sista kolumnen; för maskinklassade rader anges modellens konfidens. Databaser: OpenAlex (OA; inte open access), DBLP, Web of Science (WoS), Scopus.}\label{tab:sok-komposition}\\
\toprule
Titel & År & Leverantör & Skäl \\
\midrule
\endfirsthead
\multicolumn{4}{@{}l}{\emph{\tablename~\thetable{} (forts.)}}\\
\toprule
Titel & År & Leverantör & Skäl \\
\midrule
\endhead
\midrule \multicolumn{4}{r@{}}{\emph{forts.\ på nästa sida}}\\
\endfoot
\bottomrule
\endlastfoot
An empirical investigation of an object-oriented software system & 2000 & OA & \emph{citerad}: empiri (empirical; inheritance-defects; low-inheritance-use) \\
Encapsulation and inheritance in object-oriented programming languages & 1986 & OA & \emph{citerad}: arvets kostnader (encapsulation; inheritance-costs; change-safety) \\
Evaluating inheritance depth on the maintainability of object-oriented software. & 1996 & DBLP & \emph{citerad}: empiri (inheritance-depth; maintainability; empirical) \\
Experimental assessment of the effect of inheritance on the maintainability of object-oriented systems. & 2000 & DBLP & \emph{citerad}: empiri (inheritance-maintainability; empirical) \\
Fragile base-class problem, problem? & 2017 & Scopus & \emph{citerad}: empiri (fragile-base-class; empirical-reassessment) \\
Inheritance Usage Patterns in Open-Source Systems & 2018 & Scopus & \emph{citerad}: praktik (inheritance-usage-empirical) \\
Maintenance support for object-oriented programs & 1992 & OA & \emph{citerad}: arvets kostnader (maintenance; inheritance-costs; program-understanding) \\
A study of the fragile base class problem & 1998 & OA & stöder påståendet (0.97) (fragile-base-class; inheritance-costs) \\
A TYPE SAFE DESIGN TO ALLOW THE SEPARATION OF DIFFERENT RESPONSIBILITIES INTO PARALLEL HIERARCHIES & 2011 & WoS, Scopus & stöder påståendet (0.96) (tease-apart-inheritance; delegation; lsp) \\
Analysis of inheritance anomaly in object-oriented concurrent programming languages & 1993 & OA & stöder påståendet (0.97) (inheritance-anomaly; concurrency; inheritance-costs) \\
Bridge Pattern & 2022 & OA & stöder påståendet (0.99) (bridge-pattern; composition-over-inheritance) \\
C++ Programming with Design Patterns Revealed & 2001 & OA & stöder påståendet (0.80) (book; object-composition; reusability) \\
Combined Modelling and Programming Support for Composite States and Extensible State Machines & 2015 & WoS, Scopus & stöder påståendet (0.91) (state-pattern; delegation; composite-states) \\
Design Patterns in Object-Oriented Finite Element Programming & 2009 & OA, Scopus & stöder påståendet (0.97) (finite-element; avoid-inheritance-problems; object-composition) \\
ECS Game Engine Design & 2014 & OA & stöder påståendet (0.99) (ecs; game-engine; coupling) \\
Evolution of object behavior using context relations & 1998 & WoS, Scopus & stöder påståendet (0.97) (gof-reference; awkwardness-of-inheritance; dynamic-evolution) \\
Extending Java to dynamic object behaviors & 2003 & OA & stöder påståendet (0.90) (dynamic-behavior; decorator-inspired; compose-behaviors) \\
Extending object-oriented systems with roles & 1996 & OA & stöder påståendet (0.95) (role-hierarchies; class-hierarchy-explosion; evolving-objects) \\
Modularizing Different Responsibilities into Separate Parallel Hierarchies & 2013 & WoS, Scopus & stöder påståendet (0.96) (tease-apart-inheritance; delegation; lsp) \\
PODIO: An Event-Data-Model Toolkit for High Energy Physics Experiments & 2017 & OA & stöder påståendet (0.78) (avoid-deep-hierarchies; data-model) \\
Selected OOP Design Best Practices & 2023 & OA & stöder påståendet (0.97) (best-practices; python) \\
The Decorator Pattern in Software Engineering: Principles, Design, and Applications & 2025 & OA & stöder påståendet (0.99) (decorator; composition-over-inheritance) \\
Traits & 2006 & OA & stöder påståendet (0.96) (traits; inheritance-problems; reuse) \\
Transforming inheritance into composition. & 1999 & DBLP & stöder påståendet (0.91) (refactoring; transform-inheritance-to-composition) \\
Working with Protocols & 2011 & OA & stöder påståendet (0.98) (protocols; favoring-composition) \\
A Modified Inheritance Mechanism Enhancing Reusability and Maintainability in Object-Oriented Languages. & 1996 & DBLP & kvalificerar eller motsäger påståendet (0.83) (modified-inheritance; reusability; maintainability) \\
Blessing or Curse? Investigating Test Code Maintenance Through Inheritance and Interface & 2024 & Scopus & kvalificerar eller motsäger påståendet (0.82) (test-code-maintenance; inheritance-vs-interface) \\
Class composition for specifying framework design & 1999 & WoS, Scopus & kvalificerar eller motsäger påståendet (0.96) (frameworks; inheritance-and-composition-limits) \\
Class Inheritance Structures and Software Maintenance Activities - An Empirical Analysis & 2015 & WoS & kvalificerar eller motsäger påståendet (0.95) (inheritance-maintenance; deep-hierarchies; empirical) \\
Comparison of Multiple Inheritance in Seven Object-Oriented Programming Languages & 2026 & Scopus & kvalificerar eller motsäger påståendet (0.76) (multiple-inheritance; language-comparison) \\
Composition and inheritance are different & 2015 & Scopus & kvalificerar eller motsäger påståendet (0.78) (inheritance-vs-composition) \\
Designing with inheritance and composition. & 2012 & DBLP & kvalificerar eller motsäger påståendet (0.83) (designing-with-both; trade-offs) \\
Evaluation of design pattern alternatives in Java & 2022 & WoS & kvalificerar eller motsäger påståendet (0.96) (design-pattern-alternatives; inheritance-complexity; modern-language-features) \\
Exploring the Impact of Inheritance on Test Code Maintainability. & 2024 & DBLP & kvalificerar eller motsäger påståendet (0.83) (test-code-maintainability; inheritance) \\
Inheritance software metrics on smart contracts & 2020 & OA & kvalificerar eller motsäger påståendet (0.98) (metrics; smart-contracts; inheritance-harmless) \\
Inheritance versus Delegation: Which is more energy efficient? & 2020 & Scopus & kvalificerar eller motsäger påståendet (0.90) (inheritance-vs-delegation; energy-efficiency) \\
Leveraging Interfaces in Java to Address the Absence of Multiple Inheritance: A Comprehensive Analysis & 2025 & OA & kvalificerar eller motsäger påståendet (0.82) (multiple-inheritance; interfaces) \\
On the Evolution of Inheritance and Delegation Mechanisms and Their Impact on Code Quality & 2022 & Scopus & kvalificerar eller motsäger påståendet (0.90) (inheritance-and-delegation; code-quality-evolution) \\
Puzzle Pattern, a Systematic Approach to Multiple Behavioral Inheritance Implementation in Object-Oriented Programming & 2024 & Scopus & kvalificerar eller motsäger påståendet (0.76) (multiple-behavioral-inheritance; inheritance-benefits) \\
Reuse variables: Reusing code and state in Timor & 2004 & WoS, Scopus & kvalificerar eller motsäger påståendet (0.95) (delegation-costs; reuse; language-design) \\
Simulating multiple inheritance and generics in Java & 1999 & WoS, Scopus & kvalificerar eller motsäger påståendet (0.90) (multiple-inheritance; automatic-forwarding; trade-offs) \\
The effect of inheritance on the maintainability of object-oriented software: an empirical study. & 1995 & DBLP & kvalificerar eller motsäger påståendet (0.95) (inheritance-maintainability; empirical) \\
Toolkit design for interactive structured graphics & 2004 & OA & kvalificerar eller motsäger påståendet (0.96) (toolkit-architecture; trade-offs) \\
\end{longtable}
\endgroup

